\bookchapter{第十九章：トラブルシューティングと自己救助}

\chapterintro{一、スマホが重い、発熱——最もよく聞かれる問題}

\booksection{正直、スマホを1年ほど使うと重くなるのは正常です。原因は必ずしもスマホの故障ではなく、多くの場合「埃溜まり」——ゴミファイル、バックグラウンドプロセス、ストレージ容量が徐々に埋まる。}

まず診断、急いで修理しない：

\booksubsection{ストレージ容量を確認}

\begin{enumerate}
  \item \textbf{：設定 ストレージで残容量を確認。10\%未満だとシステムは非常に重くなります、場所を空けようと必死だから。私の経験では：15\%-20\%以上の残容量を保てば、スマホは比較的スムーズに動作。}消費電力ランキングを確認
  \item \textbf{：バッテリー設定で過去24時間の消費電力ランキングを確認。あるAppが大部分を占めているのにほとんど使っていない場合、バックグラウンドで勝手に動いています。アンインストールまたはバックグラウンド活動を制限。}最近何をしたか思い出す
  \item \textbf{：システムアップデート後に重くなった？（大型バージョン更新後の最初数日は重くなるのが正常、システムがバックグラウンドでインデックス再構築中）。それとも新しいAppをインストール後に重くなった？（アンインストールして試す）。}使用環境を確認
  \item \textbf{：スマホを太陽の下に放置、分厚いケースを付けて充電しながらゲームをすると、発熱は必然。高温はCPUの自動クロックダウンを引き起こし、スマホが突然重く感じます。冷ましてから再試行。}治療法（簡単なものから複雑なものへ）
\end{enumerate}

\textbf{：スマホを再起動（60\%の問題を解決） ストレージ容量をクリア（不要なAppと動画を削除） あるAppが問題ならキャッシュをクリアまたはアンインストールして再インストール 最後の手段：データバックアップ後に出荷時リセット。}二、アプリが落ちる、応答なし——Appが言うことを聞かない

\booksection{あるAppを開くとすぐ落ちる、または使っているうちに突然フリーズする。処理フローは医師の問診のように、一歩ずつ：}

まず閉じて再度開く

\begin{enumerate}
  \item \textbf{：バックグラウンドからAppをスワイプで消去し、再度開く。ほとんどの一時的バグはこれで解決。}Appをアップデート
  \item \textbf{：アプリストアでアップデートがあるか確認。Appバージョンが古すぎてシステムと互換性がない場合があります。}権限を確認
  \item \textbf{：一部のAppは特定の権限が必要でないと正常動作しませんが、付与していません。例えばカメラAppにカメラ権限を付与していないと開くとすぐ落ちる。設定でそのAppを見つけ、必要な権限をすべて付与。}キャッシュをクリア
  \item \textbf{：設定 アプリ管理 そのAppを見つける ストレージ キャッシュをクリア。注意：「キャッシュをクリア」であって「データをクリア」ではない！「データをクリア」はApp内のログイン状態と設定をすべて削除、Appを出荷時状態にリセットします。}アンインストールして再インストール
  \item \textbf{：最後の手段。アンインストール前にこのAppのデータが失われないか確認——WeChatなら先にチャット履歴をバックアップ、ゲームならアカウントが携帯番号に紐付いているか確認。}三、充電できない、充電が遅い——バッテリーのせいではないかも
\end{enumerate}

\booksection{充電できないのは特に焦りを引き起こす問題ですが、多くの場合原因はアクセサリにあり、スマホの故障ではありません。}

まず充電ポートを確認

\begin{itemize}
  \item \textbf{：爪楊枝や非導電の小工具で、充電ポートの埃や綿ぼこりを優しく取り出す。私自身も経験——ポケットから落ちた糸くずが充電ポートに詰まり、挿すと接触不良かと思ったが、取り出したら直った。}ケーブルと充電ヘッドを交換して試す
  \item \textbf{：ケーブルの破損率は充電ヘッドよりはるかに高い、特に頻繁に巻くケーブル。問題ないと分かっているケーブルと充電ヘッドに交換してクロステストすれば、どのアクセサリが壊れたかすぐ特定できる。}ワイヤレス充電の問題
  \item \textbf{：ワイヤレス充電を使用している場合、スマホケースが厚すぎないか、ケース裏に金属片（マグネットスタンドなど）がないか確認。また、スマホを充電パッドの中央に正確に置く——ずれると充電できない。}充電が遅い
  \item \textbf{：充電ヘッドの電力が十分か確認（iPhoneの急速充電には少なくとも20WのPD充電ヘッドが必要）、ケーブルが急速充電対応か（すべてのケーブルが3A以上の電流に対応しているわけではない）、スマホが熱すぎないか（熱すぎると自動的に出力を下げてバッテリーを保護）。}四、ネットワークに接続できない、Wi-Fiがよく切れる
\end{itemize}

\booksection{外でモバイルデータは正常だが、家に帰るとWi-Fiに接続できない。またはWi-Fiに接続しているがネットに繋がらない。この順序で確認：}

機内モードのオン/オフ

\begin{enumerate}
  \item \textbf{：コントロールセンターを下ろし、機内モードをオンにして5秒待ってからオフに。これはスマホにネットワークを再検索させるもので、非常に効果的。}Wi-Fiを「忘れる」て再接続
  \item \textbf{：Wi-Fi設定で、ネットワーク名の横の感嘆符（Androidは歯車）をタップし、「このネットワークを無視」または「忘れる」を選択、再度パスワードを入力して接続。}ルーターを再起動
  \item \textbf{：多くの家庭のルーターは数ヶ月電源を切っておらず、時間が経つと不安定に。電源を抜いて30秒待ってから再挿入、ほとんどのWi-Fi問題が解決。}スマホの問題か確認
  \item \textbf{：他のデバイスが同じWi-Fiに正常に接続できるか？他は正常で自分のスマホだけダメなら、スマホの再起動またはネットワーク設定のリセットが必要（設定 一般 転送またはリセット ネットワーク設定をリセット、データは削除されないが、すべてのWi-Fiパスワードはクリアされる）。}五、画面問題：タッチ反応不良、画面ちらつき、残像
\end{enumerate}

\booksection{タッチ反応不良}

\begin{itemize}
  \item \textbf{：まず強化ガラスフィルムを剥がして、フィルムのせいか確認——安い強化ガラスフィルムは長く使うとタッチに影響。画面に水や手に油があるとタッチがずれることも。どちらでもなければ、再起動して再試行。}画面に跡/残像（焼き付き）
  \item \textbf{：長時間同じ画面を表示（例えばずっと画面を点灯してDouyinを見る）と、OLED画面に「残像」が残ります。これは画面故障ではなく、一時的。画面を消して30分休めば通常回復。}画面ちらつき
  \item \textbf{：低輝度でちらつく場合、PWM調光の問題の可能性が高く、OLED画面で多い。設定で「DC調光」や「フリッカー防止」オプションを探し、オンにするとかなり改善。すべてのスマホにある機能ではありません。}画面の一部が完全に黒/白
  \item \textbf{：これは比較的深刻で、画面の物理的損傷の可能性が高く、画面交換が必要。非公式修理の画面は、色と輝度が通常純正に劣り、防水にも影響。}六、起動できない、Logo画面で止まる——まず焦らない
\end{itemize}

\booksection{この状況は最も心拍数を上げますが、冷静に、まずこれらを試しましょう：}

まず30分充電してから

\begin{itemize}
  \item \textbf{：多くの場合、単にバッテリーが完全に切れているだけ。純正充電器を挿し、少なくとも30分待ってから起動を試みる。一部のスマホはバッテリー切れ後、一定量まで充電されないと反応しません。}強制再起動
  \item \textbf{：全面屏iPhoneは音量上げ 音量下げ 電源キー長押しでLogo表示まで。Androidは電源キー10〜15秒長押し、振動まで。}起動画面でループ再起動
  \item \textbf{：この状況は「白リンゴ」（iPhone）または「Logo卡まり」（Android）と呼ばれます。強制再起動が効かない場合、リカバリーモードに入ってパソコンで修復が必要。iPhoneをパソコンに接続しiTunes/Finderを開き、リカバリーモードに入った後「アップデート」を選択（「復元」ではない）、システム再インストールでデータ保持。AndroidはFastbootまたはRecoveryモードへ、各ブランドでキーの組み合わせが異なるので、「スマホ型番＋リカバリーモード入り方」を検索推奨。}上記すべて効かない
  \item \textbf{：専門修理に頼るしかありません。データをバックアップしていれば、あまり心配不要。}七、音声問題：スピーカーが鳴らない、通話音が小さい
\end{itemize}

\booksection{スピーカーから音が出ない}

\begin{itemize}
  \item \textbf{：まずマナーモード（iPhone側面のサイレントスイッチ）、おやすみモード、音量が最低になっていないか確認。馬鹿げているように聞こえるが、この状況が40\%を占めます。次にBluetoothがイヤホンやスピーカーに接続していないか確認。}通話相手に声が聞こえない
  \item \textbf{：マイク穴が塞がれていないか確認——スマホ底部の小穴は埃で塞がれやすい。柔らかいブラシで優しく掃除。針や爪楊枝で強く突かない、防塵ネットを破損。}音が割れる/ノイズ
  \item \textbf{：スピーカーの浸水が原因の場合——ドライヤーの熱風を当てない！ティッシュで表面の水分を吸い取り、自然乾燥。ネットでは音波排水App（低周波音を再生して水を振動で飛ばす）がある程度効果的と言われており、試す価値あり。}八、公式修理を探す、落とし穴を避ける
\end{itemize}

\booksection{上記の方法をすべて試してもダメなら、専門修理に頼る必要があります。しかし修理市場は水深が深く、私自身の失敗経験と教訓を：}

本物の公式を見つける、「公式っぽい」ではなく

\begin{itemize}
  \item \textbf{：ブランド公式サイト（Apple.com.cn、Huawei.com.cnなど）で「認定サービスプロバイダー」リストを探す。街角でAppleロゴを掲げる非認定店は、純正アクセサリを使っていない可能性があり、修理後に防水がなくなります。}修理に出す前に必ずやること
  \item \textbf{：すべてのデータをバックアップ（修理中にデータが失われないことを期待しない）、すべての重要なアカウントをログアウト（Apple ID、WeChat、Alipay）、「iPhoneを探す」をオフ。一部の修理店はロック画面パスワードの提供を要求しますが、プライバシーリスクあり——自分で先にアカウントをログアウトしてから持ち込むのが良い。}画面交換の落とし穴
  \item \textbf{：非公式画面は「原圧画面」（純正ガラス交換）、「後圧画面」（非純正ガラス）、「組立画面」（全非純正）に分かれる。価格は数倍違うが、表示効果と耐久性も大きく異なる。スマホがまだ保証期間内なら、大人しく公式ルートを。}バッテリー交換の落とし穴
  \item \textbf{：非純正バッテリーは容量を偽装している可能性があり、半年で純正バッテリーの2年分の劣化を示す。Appleバッテリー交換は公式価格が非公式よりそれほど高くなく、公式推奨。}データ復元の落とし穴
  \item \textbf{：スマホの浸水/落下破損後、修理店が「データ復元は追加料金」と言う場合、警戒——一部の悪徳業者はわざと損傷程度を誇張して多く請求。正規のデータ復元は先に見積もりしてから操作、復元失敗は無料。}万能の第一の手段：スマホを再起動
\end{itemize}

\begin{tcbwarn}{どんな問題に直面しても——動作が重い、アプリが落ちる、ネットワークに繋がらない、音声異常——まずスマホを再起動。これは適当な対応に聞こえるが、私自身の経験では、少なくとも6割の小さな問題は再起動で解決。スマホもパソコンと同じで、長く動作するとメモリに様々なゴミが蓄積し、再起動はクリアしてやり直すこと。だから次にスマホに問題が起きたら、急いで「どうしよう」と検索する前に、再起動してから確認を。}
今すぐ覚える：これらの救命操作は完全に暗記すべき

\end{tcbwarn}

\begin{tcbpractice}{強制再起動}
\begin{enumerate}
  \item \textbf{：全面屏iPhoneは音量上げ 音量下げ 電源キー長押しでLogo表示まで。Androidは電源キー10〜15秒長押し、振動まで。}充電ポートを掃除
  \item \textbf{：小工具（爪楊枝やSIM取り出しピン）で、充電ポートの埃を優しく掃除。}Wi-Fi再接続
  \item \textbf{：Wi-Fi設定で現在のWi-Fiを「忘れる」、再度パスワードを入力して接続。}機内モードのオン/オフ
  \item \textbf{：コントロールセンターを下ろし、機内モードをオン、心の中で5秒数え、オフにしてネットワークが復帰するか確認。}これらの操作手順をスクショで保存するか、メモに書く。問題発生時に慌てて検索する必要がなく、直接メモを見ればOK。
  \item 第二十章：開発者モードと高度な権限
\end{enumerate}

\end{tcbpractice}
